% TikZ 图：光电效应
% Day 27：近代物理
\begin{tikzpicture}[scale=0.9,>=Latex]
  \draw[fill=gray!20,thick] (0,0) rectangle (3.8,0.35);
  \node at (1.9,-0.3) {金属板};
  \foreach \x in {0.6,1.25,1.9}{
    \draw[->,blue,thick] (\x,2.25) -- (\x+0.4,0.45);
    \node[blue] at (\x-0.1,2.45) {$h\nu$};
  }
  \foreach \x/\ang in {1.0/55,1.7/75,2.4/42}{
    \draw[->,red,thick] (\x,0.36) -- ++({1.0*cos(\ang)},{1.0*sin(\ang)});
    \fill[red] (\x,0.36) circle (1.6pt);
  }
  \node[red] at (3.75,1.45) {光电子};
  \node[align=left,font=\small] at (6.25,1.1)
    {$h\nu=W_0+E_{k\max}$\\[4pt]
     $\nu\ge\nu_0$ 时\\能发生光电效应。};
\end{tikzpicture}
